%*****************************************
\chapter{Related Work}
\label{ch:relatedwork}
%*****************************************
This chapter gives a comprehensive overview of the related work by other authors
and analyzes why the existing approaches are not capable of solving the problem
defined in the introduction, thereby defining the research gap. The review is
organized along four complementary lines of research that jointly span the
design space of the thesis:
\begin{itemize}
    \item \textbf{Remote physiological sensing} (Section~\ref{sec:rw-sensing})
    traces the evolution of remote physiological sensing and establishes the
    modalities that the thesis seeks to fuse.
    \item \textbf{Masked Autoencoders for spatiotemporal data}
    (Section~\ref{sec:rw-mae}) reviews Masked Autoencoders (MAEs), the
    architectural paradigm underlying the proposed approach.
    \item \textbf{Multimodal learning and joint embedding spaces}
    (Section~\ref{sec:rw-multimodal}) examines how information from
    heterogeneous sensors can be shared across modalities.
    \item \textbf{Multiscale spectral modeling} (Section~\ref{sec:rw-spectral})
    complements this discussion with techniques that explicitly exploit the
    structure of physiological signals.
\end{itemize}
Finally, Section~\ref{sec:rw-analysis} synthesizes these threads and defines
the research gap addressed by this thesis.

\section{Evolution of Remote Physiological Sensing}
\label{sec:rw-sensing}
The field of remote photoplethysmography (rPPG) has evolved from handcrafted
mathematical models to deep neural architectures. Early supervised models such
as PhySU-Net demonstrate that vital signs can be extracted directly
from facial video \cite{savic2024physunet}; however, these models often lack
interpretability and degrade under varying environmental conditions. To address
these limitations, PHASE-Net introduces a physics-informed design that grounds
its attention mechanism in the Navier--Stokes equations, modeling pulse dynamics
as second-order damped harmonic oscillators \cite{zhao2024phasenet}.

Beyond algorithmic advances, the hardware used for physiological sensing has also
evolved: because frame-based sensors sample at only $30$--$60$ frames per
second, they impose a fundamental limit on the temporal resolution of the
captured signal. Recent work based on the EMPD dataset therefore introduces
event-based neuromorphic sensing, which offers microsecond-level temporal
resolution and bridges this sampling gap at the hardware level
\cite{feng2026empd}.

Neither approach can compensate for a sensor that has failed completely. These
limitations correspond exactly to the scenario investigated in this thesis and
highlight the need for \emph{cross-modal} redundancy rather than sensor-level
redundancy alone.

\section{Masked Autoencoders for Spatiotemporal Data}
\label{sec:rw-mae}
The success of self-supervised learning (SSL) in natural language processing has
been ported to the visual domain through Masked Autoencoders (MAEs)
\cite{he2022masked}. Whereas the original MAE randomly masks a large fraction of
spatial patches in still images, VideoMAE extends this paradigm to video and
shows that the high temporal redundancy of video data makes extremely high
masking ratios ($90\%$--$95\%$) not only viable but beneficial: aggressive
masking forces the model to reason about high-level spatiotemporal structure
rather than relying on local temporal smoothness \cite{tong2022videomae}. VideoMAE masks entire \emph{spatiotemporal tubes} that span the
temporal dimension, so that the reconstruction task cannot be solved by copying
information across adjacent frames. This tube-masking strategy is directly
relevant to the temporal masking investigated in Sub-RQ~2 of this thesis. 

VideoMAE V2 later scales this framework and introduces a dual-masking strategy that masks both encoder
and decoder inputs, substantially reducing the computational overhead of
billion-parameter models \cite{wang2023videomaev2}. 

For multi-sensor environments, MV2MAE extends the MAE paradigm to synchronized multi-view videos,
using a cross-view reconstruction task to inject geometric information and to
ensure representation robustness to viewpoint shifts \cite{shah2023mv2mae}.

The common thread of these works is that the input and target modalities
coincide: the model reconstructs the same stream from which the visible context
was drawn. They therefore establish strong within-modality priors, but leave the
question of \emph{cross-modal} reconstruction--recovering a completely missing
stream from the remaining sensors--largely unaddressed.

\section{Multimodal Learning and Joint Embedding Spaces}
\label{sec:rw-multimodal}
This thesis is situated at the intersection of heterogeneous modalities (RGB, 3D
geometry, and thermal infrared), so multimodal learning frameworks are of
central relevance. MultiMAE shows that pre-training a single Transformer on
multiple dense visual tasks (e.g., RGB, depth, and semantic segmentation) yields
a generalist model that remains robust to unpredicted modality dropout at
inference time \cite{bachmann2022multimae}. 

ImageBind provides a complementary framework for learning joint embeddings across six disparate
modalities---including thermal and IMU---by using images as a \enquote{binding}
modality that aligns signals that are rarely observed together
\cite{girdhar2023imagebind}. 

Moving beyond symmetric alignment, xMAE reframes cross-modal learning as a problem of \emph{asymmetric temporal observability}:
one signal (e.g., PPG) is used as training-time scaffolding to reconstruct a
directionally dependent second signal (e.g., ECG), thereby capturing the
electrical-to-mechanical delay inherent in human physiology \cite{zhou2026xmae}.

These frameworks establish that shared embedding spaces can align physiologically related modalities and that a model can tolerate the absence of
one modality during inference. However, they treat missingness as a robustness
problem: the model is expected to produce a \emph{usable representation} despite
dropout, not to \emph{generate a faithful reconstruction} of the missing raw
stream. The fidelity of the recovered signal is therefore neither optimized nor
quantified, which lies at the heart of the main research question (RQ) of this
thesis.

\section{Multiscale Spectral Modeling of Physiological Signals}
\label{sec:rw-spectral}
Because physiological signals are non-stationary and quasi-periodic, standard
time-domain masking is often insufficient. bioFAME addresses this by introducing
frequency-aware Transformers that parameterize representations in the frequency
domain, ensuring robustness to modality mismatch in wearable data
\cite{liu2024biofame}. The MAEConformer architecture combines convolutional
local feature extraction with Transformer-based long-range attention and employs
a Multi-Resolution STFT (MR-STFT) loss to learn temporal and spectral
characteristics jointly \cite{yu2026maeconformer}. 

To model the hierarchical rhythms of biosignals explicitly, the MMR (Wavelet-Driven Masked Multiscale
Reconstruction) framework applies the Discrete Wavelet Transform (DWT); by
reconstructing masked wavelet coefficients, MMR captures features ranging from
fine-grained waveform morphology to global rhythmic dynamics, providing a robust
foundation for downstream health monitoring \cite{thukral2026wavelet}. Finally,
rPPG-MAE adapts these self-supervised paradigms to remote measurement by
leveraging the self-similar priors inherent in biological pulse signals
\cite{liu2023rppgmae}.

These methods refine the reconstruction objective with spectral losses and show
that exploiting the quasi-periodic structure of physiological signals improves
reconstruction fidelity. Nevertheless, they operate almost exclusively within a
single modality and on wearable or contact-based signals. None of them addresses
the cross-modal setting in which a high-dimensional visual stream (e.g., thermal
infrared) must be generated from structurally different surviving modalities
(e.g., RGB video), nor do they quantify the physiological information lost in
such a recovery.

%%
%%
\section{Analysis of Related Work}
\label{sec:rw-analysis}
The existing literature demonstrates that MAEs are highly effective at learning
spatiotemporal representations for both visual \cite{tong2022videomae,
wang2023videomaev2} and physiological data \cite{liu2023rppgmae, savic2024physunet},
and that multimodal joint embedding spaces can align heterogeneous physiological
modalities \cite{bachmann2022multimae, girdhar2023imagebind, zhou2026xmae}.
However, current state-of-the-art approaches are predominantly designed for
representation learning---in which the reconstruction task serves as a pretext
to improve downstream classification or regression accuracy---rather than for
the explicit generative recovery of a missing sensor stream. This leaves three
critical gaps that prevent existing models from addressing the main research
question (RQ) regarding targeted sensor failure and information loss:

\begin{itemize}
    \item \textbf{The Downstream vs.\ Generative Objective (RQ):} Most
        physiological MAEs, such as MAEConformer \cite{yu2026maeconformer} and
        bioFAME \cite{liu2024biofame}, evaluate success by performance on
        specific medical tasks, for example HIE severity classification or
        cardiovascular outcome prediction. While xMAE \cite{zhou2026xmae} uses
        cross-modal reconstruction as a scaffolding task, its primary goal is to
        capture directional inductive biases that enhance the PPG representation,
        not to quantify the fidelity of the reconstructed signal for its own
        sake. There is a lack of research that treats the generative
        reconstruction of the raw signal as the primary metric, i.e., that
        quantifies how much biological information (e.g., HRV morphology,
        hemodynamic transients) can be recovered when a physical sensor fails.
        This gap maps directly onto the main RQ.

    \item \textbf{Robustness vs.\ Active Reconstruction during Failure (RQ, Sub-RQ~1):}
        Multimodal frameworks such as MultiMAE \cite{bachmann2022multimae} and
        ImageBind \cite{girdhar2023imagebind} focus on binding disparate
        modalities into a joint embedding space to achieve robustness against
        unpredicted modality dropout. These models continue to function if a
        sensor fails, but they are not optimized for cross-modal predictive
        coding---specifically, the use of surviving modalities (e.g., RGB video
        and 3D geometry) to actively reconstruct a missing high-dimensional
        sensor stream (e.g., thermal infrared). Existing work treats missingness
        as a noise-insensitivity problem, whereas the RQ of this thesis frames
        it as a structured inference problem that must be solved by a single
        unified encoder capable of fusing 1D time series with 2D video
        (Sub-RQ~1).

    \item \textbf{Lack of Information-Loss Quantification in the Signal Domain (RQ):}
        The research community currently lacks a standardized framework for
        measuring information loss during sensor failure in non-invasive
        sensing. While MMR \cite{thukral2026wavelet} and MAEConformer
        \cite{yu2026maeconformer} use spectral losses (DWT and STFT) to improve
        reconstruction, they do not quantify how much physiological signal
        information is systematically discarded when a modality must be predicted
        across a spectral gap (e.g., reconstructing thermal heat signatures from
        RGB color changes). Existing metrics focus on error (MAE/RMSE) or
        correlation \cite{savic2024physunet}, but they do not address the entropy
        or physiological completeness of the reconstructed stream during a
        targeted hardware failure.
\end{itemize}

\subsection{Summary and Research Gap}
While the community has explored within-modality interpolation and multimodal
task performance, there is a distinct gap in understanding cross-modal
predictive limits. This thesis fills this gap by shifting the focus from
\enquote{representations for classification} to \enquote{predictive coding for
sensor recovery}. Concretely, the thesis contributes:
\begin{itemize}
    \item a unified Vision Transformer encoder that fuses 1D physiological
    time series (e.g., EDA, heart rate) with 2D RGB and thermal video
    (Sub-RQ~1);
    \item a systematic study of temporal tube masking and extreme masking
    ratios ($90\%$ and above) for the reconstruction of physiological
    indicators from noisy video (Sub-RQ~2); and
    \item a first rigorous analysis, based on simulated sensor failures in the
    BP4D+ corpus \cite{zhang2014bp4d}, of the extent to which the cross-modal
    redundancy of human physiology enables a deep model to compensate for
    missing hardware.
\end{itemize}

\begin{table}[htbp]
\centering
\small
\setlength{\tabcolsep}{5pt}
\caption{Comprehensive overview of the surveyed methods. \enquote{Within/Cross-modal}
indicates whether the reconstruction target is drawn from the same modality as the
visible context or from a different one; \enquote{Sensor-failure focus} states whether a
method explicitly handles a completely missing sensor stream. The final row
positions this thesis relative to the state of the art.}
\label{tab:rw-overview}
\begin{tabular}{@{}p{3.2cm} p{2.5cm} p{4.3cm} p{2.2cm} p{2.4cm}@{}}
\toprule
Method & Modalities & Reconstruction objective & Within/Cross-modal & Sensor-failure focus \\
\midrule
PhySU-Net \cite{savic2024physunet}
    & RGB face video & supervised rPPG estimation & Within & No \\
PHASE-Net \cite{zhao2024phasenet}
    & RGB face video & physics-informed rPPG estimation & Within & No \\
EMPD \cite{feng2026empd}
    & Event camera & event-driven rPPG estimation & Within & No \\
VideoMAE \cite{tong2022videomae}
    & Video & tube-masked reconstruction (90--95\,\%) & Within & No \\

MV2MAE \cite{shah2023mv2mae}
    & Multi-view video & cross-view reconstruction & Cross-view & Partial \newline (viewpoint) \\
MultiMAE \cite{bachmann2022multimae}
    & RGB, depth, segmentation & multi-task masked reconstruction & Cross-modal & Partial (dropout) \\
ImageBind \cite{girdhar2023imagebind}
    & Six modalities (incl.\ thermal, IMU) & joint embedding alignment & Cross-modal & Partial \newline (no generation) \\
xMAE \cite{zhou2026xmae}
    & PPG, ECG & asymmetric cross-modal scaffolding & Cross-modal & No \\
bioFAME \cite{liu2024biofame}
    & Wearable biosignals & frequency-aware reconstruction & Within & No \\
MAEConformer \cite{yu2026maeconformer}
    & Biosignals & MR-STFT masked reconstruction & Within & No \\
MMR \cite{thukral2026wavelet}
    & Biosignals & DWT coefficient reconstruction & Within & No \\
rPPG-MAE \cite{liu2023rppgmae}
    & Face video / rPPG & self-similar masked reconstruction & Within & No \\
\midrule
\textbf{This thesis}
    & \textbf{RGB, 3D, thermal video, 1D series}
    & \textbf{generative cross-modal recovery}
    & \textbf{Cross-modal} & \textbf{Yes (targeted failure)} \\
\bottomrule
\end{tabular}
\end{table}

Table~\ref{tab:rw-overview} condenses the reviewed literature along four
dimensions that are decisive for the problem of this thesis: the modality of the
input, the reconstruction objective, whether the reconstruction target lies
within or across modalities, and whether a completely missing sensor stream is
explicitly handled. Two patterns emerge. First, the majority of methods operate
\emph{within} a single modality---reconstructing masked parts of the same
stream, be it video (VideoMAE, MV2MAE) or physiological signal
(bioFAME, MAEConformer, MMR, rPPG-MAE)---and therefore cannot recover a stream
that is entirely absent. Second, the methods that do operate \emph{across}
modalities (MultiMAE, ImageBind, xMAE) treat missingness either as a robustness
requirement during pre-training or as a scaffolding objective for representation
learning; none of them optimizes or quantifies the generative fidelity of a
reconstructed sensor stream. As a result, no existing approach combines a unified
encoder for 1D time series and 2D video with an explicit, measurable generative
recovery of a failed sensor---the combination this thesis investigates.